\section{Expected Contributions}
\label{sec:contributions}

At the end of this work the following will exist, none of which exists now.

\begin{enumerate}[leftmargin=*,itemsep=5pt]
  \item \textbf{The first systematic map of harm communication in health
    tracking apps.} A coded account, across 50 fitness and fertility tracking
    apps in the German market, of which privacy harm types providers name, which
    they minimise, and which never appear --- broken down by harm category, by
    domain, and by communication surface. This answers the question the harm
    literature and the tracker-measurement literature both leave open.

  \item \textbf{A quantified account of four specific silences.} Per-app and
    per-surface figures for whether inference, breach consequences, incident
    handling, and any concrete consequence for the user are mentioned at all.
    This is the empirical core: it converts the observation motivating the
    thesis into a number that can be cited, contested, and re-measured later.

  \item \textbf{A reusable harm-communication codebook.} A versioned coding
    instrument that operationalises the Citron and Solove
    typology~\parencite{citron2021harms} for corporate privacy text --- each
    category with a definition, a positive example, and a counter-example ---
    extended with the Knabenhans et al.\ pathway
    annotation~\parencite{knabenhans2026harms} and a PET-framing axis derived
    from Martin et al.~\parencite{martin2025cookies}. Existing annotation
    schemes for privacy policies code data practices; this one codes harm, and
    is designed to be applied to other domains unchanged.

  \item \textbf{A privacy-focused cognitive walkthrough instrument.} An
    adaptation of Wharton et al.'s
    procedure~\parencite{wharton1994cognitive} in which the user goals are
    privacy goals --- find out what is inferred, find out what happens after a
    breach, stop third-party sharing, delete everything --- with a fixed task
    sequence and record sheet. It is reusable by anyone inspecting privacy
    communication in an interface rather than in a document.

  \item \textbf{A reproducible LLM-assisted analysis pipeline, honestly
    reported.} Versioned prompts, a pinned model, structured output with
    mandatory verbatim spans, a run log, and a validation report giving Cohen's
    $\kappa$ per category against manual coding together with the run-to-run
    disagreement rate and the precision audit. Whatever the agreement figures
    turn out to be, they are published with the results --- which is what makes
    the pipeline usable by the next person rather than only by its author.

  \item \textbf{Design and policy implications.} Concrete statements about where
    harm-relevant information is missing from the surface on which users
    actually decide, and what a disclosure regime built around data practices
    does not cause providers to say. If a second phase is pursued, the coded
    corpus supplies realistic, evidence-based stimulus material for the
    scenario-based survey described in the original thesis offer.
\end{enumerate}
